\documentclass[10pt,a4paper]{amsart}
\usepackage[margin=19mm]{geometry}
\usepackage{amsmath,amssymb,mathtools}
\usepackage[scr=rsfs,cal=euler]{mathalfa}
\usepackage{xfrac}
\usepackage[unicode,hidelinks,bookmarks=false]{hyperref}
\newcommand{\ie}{i.e.\ }
\newcommand{\cf}{cf.\ }
\newcommand{\resp}{respectively\ }
\newcommand{\Subsection}{\subsection}
\providecommand{\nobreakdash}{\nobreak\hbox{-}}
\setlength{\emergencystretch}{2em}
\multlinegap0pt
\usepackage{amssymb}
\usepackage{float}
\usepackage{dsfont}
\newtheorem{thm}{Theorem}
\newtheorem{lem}{Lemma}
\newtheorem{pro}{Proposition}
\title{\bf The range of once-reinforced random walk on the half-line}
\author{Zechun Hu, Ting Ma, Renming Song, Li Wang}
\date{Source: arXiv:2602.10721v2 (CC BY 4.0)}
\begin{document}
\maketitle

\noindent\emph{Source and licence.} \bf The range of once-reinforced random walk on the half-line. Version arXiv:2602.10721v2 (CC BY 4.0), distributed by arXiv under the Creative Commons Attribution 4.0 International licence, http://creativecommons.org/licenses/by/4.0/, as recorded in the arXiv licence metadata for this exact version and shown on the abstract page https://arxiv.org/abs/2602.10721v2. Full licence notice: \texttt{licenses/arxiv-2602.10721-CC-BY-4.0.txt}.\par\smallskip

\noindent\textit{Editorial scope.} Counted unit: the article's thm thm-1.1 (source0.tex lines 148--162). The article's complete original proof of it is delivered with it (source0.tex lines 389--412, one \texttt{proof} environment of 24 lines, which the article places away from the statement). No mathematics is written, repaired or re-derived: the statement and the proof are the article's own lines, in the article's own notation, and no retained line is substituted or re-flowed.  Retained uncounted as the source-local context the proof needs: lines 96--98; lines 224--237; lines 234--236; lines 366--372; lines 376--385; lines 382--384.  Nothing was omitted from any retained span, so the package carries no omission record.  No retained line contains a citation, so the wrapper carries no bibliography and no bibliography entry.  No retained line contains a figure environment, an image-inclusion command or drawing code, so no image is delivered and the package declares no asset dependency.  Editorial formatting only, in addition: an AMS \texttt{amsart} preamble that restates the article's own declarations and macros used by the retained lines, namely the environments \texttt{thm}, \texttt{lem}, \texttt{pro} on separate counters, together with the standard packages those lines need.  The retained environments print in this document as Theorem 1, Lemma 1, Pro 1 in order of appearance, whereas the article prints the same items with the numbering of its own full document.  The complete native source of the article as distributed in the arXiv e-print for this exact version is bundled unchanged as \texttt{sources/194422/source0.tex}.
\begin{equation}\label{2.2}
R_{n}:= \max\{X_{0},X_{1},\ldots, X_{n}\}.
\end{equation}

\begin{thm}\label{thm-3.1}
Let $a_{1}, a_{2}, \ldots \geq 0$ such that $\sum_{n=1}^{\infty} a_{n} x^{n}$ converges for $|x|<1$. Suppose that for some $\alpha, A \geq 0$,
\begin{equation*}
\sum_{n=1}^{\infty} a_{n} x^{n} \stackrel{x \uparrow 1}{\sim} \frac{A}{(1-x)^{\alpha}}.
\end{equation*}
Then,
\begin{equation*}
\sum_{k=1}^{n} a_{k} \stackrel{n \rightarrow \infty}{\sim} \frac{A}{\Gamma(\alpha+1)} n^{\alpha}.
\end{equation*}
Moreover, if $\alpha>1$, and $n \mapsto a_{n}$ is nondecreasing,
\begin{equation}\label{3.2-a}
a_{n} \sim \frac{A \alpha}{\Gamma(\alpha+1)} n^{\alpha-1}.
\end{equation}
\end{thm}

\begin{equation}
a_{n} \sim \frac{A \alpha}{\Gamma(\alpha+1)} n^{\alpha-1}.
\end{equation}

\begin{lem}\label{lem-3.4}
For $s \in(0,1)$ and $ \ell=0,1,2, \ldots$, it holds
that
\begin{equation}\label{3.15-a}
H_{\ell}(s):=\sum_{n=1}^{\infty} s^{n} \mathbb{E}\left[R_{n} \cdots\left(R_{n}+\ell\right)\right]=\frac{\ell+1}{1-s} \sum_{k=1}^{\infty} k \cdots(k+\ell-1) s \prod_{i=1}^{k-1} G_{i}(s).
\end{equation}
\end{lem}

\begin{pro}\label{pro-3.5}
Let $H_\ell$ for $\ell = 0, 1, 2, \ldots$ be as defined in Lemma \ref{lem-3.4}. Then,
\begin{equation*}
H_\ell(s) \stackrel{s \uparrow 1}{\sim} K_{\ell}(1-s)^{-(\ell+3)/2} ,
\end{equation*}
where
\begin{equation}\label{3.15}
K_{\ell}=\frac{\ell+1}{2^{(\ell+1)/2}}2^{c}\int_{0}^{\infty}x^{\ell}\left(\frac{e^{x}}{e^{2x}+1}\right)^{c}dx.
\end{equation}
\end{pro}

\begin{equation}
K_{\ell}=\frac{\ell+1}{2^{(\ell+1)/2}}2^{c}\int_{0}^{\infty}x^{\ell}\left(\frac{e^{x}}{e^{2x}+1}\right)^{c}dx.
\end{equation}

\begin{thm}\label{thm-1.1}
 Let $R_{n}$, defined in \eqref{2.2}, be  the range of the ORRW on $\mathbb{Z}_{+}$ with parameter $c>0$. Then,
\begin{equation*}
\mathbb{E}\left[\left(\frac{R_{n}}{\sqrt{n}}\right)^{\ell}\right] \sim \frac{1}{2^{(\ell-2)/2}
\Gamma(\ell/2)} \cdot J_{\ell}(c),\quad \ell=1,2, \ldots,
\end{equation*}
where
\begin{equation*}
J_{\ell}(c):=2^{c} \int_{0}^{\infty} x^{\ell-1}\left(\frac{e^{x}}{e^{2x}+1}\right)^{c}dx.
\end{equation*}
In particular,
\begin{equation*}
\mathbb{E}\left[\frac{R_{n}}{\sqrt{n}}\right] \sim \sqrt{\frac{2}{\pi}} J_{1}(c) \quad \text { and } \quad \mathbb{E}\left[\frac{R_{n}^{2}}{n}\right] \sim J_{2}(c).
\end{equation*}
\end{thm}

\begin{proof}[Proof of Theorem \ref{thm-1.1}]

Let $a_n := \mathbb{E}[R_n \cdots (R_n + \ell)]$. It follows from Proposition \ref{pro-3.5} that
\begin{equation*}
H_\ell(s) = \sum_{n=1}^\infty a_n s^n \stackrel{s \uparrow 1}{\sim} \frac{K_\ell}{(1-s)^{(\ell+3)/2}},
\end{equation*}
Thus, Theorem \ref{thm-3.1} holds with $A=K_{\ell}$ and $\alpha = (\ell+3)/2$. In particular, \eqref{3.2-a} implies that
\begin{equation*}
\mathbb{E}[R_n \cdots (R_n + \ell)]\sim
\frac{K_{\ell}\cdot (\ell +3)/2}{\Gamma((\ell+5)/2)}n^{(\ell+1)/2}.
\end{equation*}
Combining \eqref{3.15} with the identity  $\Gamma(x+1)=x\Gamma(x)$, the coefficient
can be simplified  to
\[
\frac{K_{\ell}\cdot (\ell +3)/2}{\Gamma((\ell+5)/2)}
=\frac{2^{c}}{2^{(\ell-1)/2}\Gamma((\ell+1)/2)} \int_{0}^{\infty}x^{\ell }\left( \frac{e^{x}}{e^{2x}+1} \right)^{c}dx.
\]
Consequently,  for  every integer  $\ell=0, 1, 2,\ldots$,
\[
\mathbb{E}[R_{n}^{\ell+1}] \sim
\frac{2^{c}}{2^{(\ell-1)/2}\Gamma((\ell+1)/2)} \int_{0}^{\infty}x^{\ell }\left( \frac{e^{x}}{e^{2x}+1} \right)^{c}dx  \cdot n^{(\ell+1)/2}.
\]
The proof is complete.
\end{proof}

\end{document}
